%==============================================================================
%  ANNEXE - Choix technologique : conformite academique et vision marche
%  Migration J2EE -> Spring Boot (tableau de correspondance + argumentaire)
%  A inserer via %==============================================================================
%  ANNEXE - Choix technologique : conformite academique et vision marche
%  Migration J2EE -> Spring Boot (tableau de correspondance + argumentaire)
%  A inserer via %==============================================================================
%  ANNEXE - Choix technologique : conformite academique et vision marche
%  Migration J2EE -> Spring Boot (tableau de correspondance + argumentaire)
%  A inserer via %==============================================================================
%  ANNEXE - Choix technologique : conformite academique et vision marche
%  Migration J2EE -> Spring Boot (tableau de correspondance + argumentaire)
%  A inserer via \input{annexe_migration_springboot} dans cahier_des_charges.tex
%  (utilise la charte "Miel & noir" : miel, mielfonce, encadre, colonne C,
%   longtable, booktabs -- deja definis dans le preambule principal)
%==============================================================================

%==============================================================================
\chapter{Annexe - Choix technologique : academique et vision marche}
\label{chap:migration}
%==============================================================================
Cette annexe repond a la question du \emph{pourquoi} et du \emph{comment} du
langage de developpement. Elle confronte le socle \textbf{J2EE} impose par le
present cahier des charges a l'ecosysteme \textbf{Spring Boot}, devenu le
standard du developpement Java d'entreprise, et documente une trajectoire de
modernisation sans changement de langage.

\section{Positionnement}
Le systeme \textbf{Z\"umm} est implemente en \textbf{J2EE (Java EE)}
conformement aux exigences techniques (chapitre~\ref{chap:conception} et grille
d'evaluation) : MVC, Servlets, JSP, EJB, JDBC, internationalisation
externalisee, services web securises (OAuth, X.509 / SSL) et API interrogeable
par des applications tierces. Ce socle materialise les competences fondamentales
du parcours DSI : maitrise des couches d'un systeme d'information, separation des
responsabilites et interoperabilite.

\begin{encadre}
\textbf{Constat marche :} les briques Java EE historiques (JSP, Servlets
\og nus \fg{} et EJB) sont aujourd'hui considerees comme \emph{heritees}
(\emph{legacy}). Les recrutements sur les nouveaux systemes d'information
privilegient tres majoritairement l'ecosysteme \textbf{Spring Boot}. La presente
annexe conserve le langage \textbf{Java} et demontre, techno par techno,
l'equivalent moderne de chaque exigence.
\end{encadre}

%------------------------------------------------------------------------------
\section{Tableau de correspondance : coeur technique}
%------------------------------------------------------------------------------
Le tableau suivant met en regard chaque technologie imposee (et notee) avec son
equivalent Spring Boot, ainsi que la nature concrete du changement.

\begin{longtable}{>{\bfseries}p{3cm} p{3.4cm} p{7.2cm}}
\toprule
\textbf{J2EE (impose)} & \textbf{Spring Boot} & \textbf{Ce qui change concretement} \\
\midrule
\endhead
MVC (manuel) & Spring MVC (\texttt{@Controller}, \texttt{@RequestMapping}) & Le pattern reste ; le routage devient declaratif par annotations au lieu du \texttt{web.xml}. \\
Servlet (\texttt{HttpServlet}, \texttt{doGet/doPost}) & \texttt{@RestController} / \texttt{@Controller} & Plus de \texttt{extends HttpServlet} : une methode annotee = un endpoint. Le \texttt{DispatcherServlet} orchestre en arriere-plan. \\
JSP (presentation) & Thymeleaf (server-side) ou React / Angular (SPA) & Templating moderne sans scriptlets Java ; pour un effet marche fort, API REST + front decouple. \\
EJB (\texttt{@Stateless}, conteneur EJB) & Spring \texttt{@Service} / \texttt{@Component} & Meme role (logique metier, injection, transactions) sans conteneur EJB lourd ; \texttt{@Transactional} gere les transactions. \\
JDBC (SQL manuel, \texttt{ResultSet}) & Spring Data JPA / Hibernate & Le mapping objet-relationnel remplace le SQL a la main ; \texttt{JpaRepository} genere les CRUD (\texttt{JdbcTemplate} reste disponible). \\
i18n par fichier XML & \texttt{messages\_fr/en/ar.properties} + \texttt{MessageSource} & Meme principe d'externalisation des textes, format \texttt{.properties} standard, toujours FR / EN / AR. \\
\texttt{ConfigZumm.ini} & \texttt{application.yml} + \texttt{@ConfigurationProperties} & Parametrage sans recompilation, injectable par profil ; les seuils et unites deviennent des beans types. \\
OAuth Google (WS XML maison) & Spring Security OAuth2 Client & Login Google en quelques lignes de configuration au lieu d'un client OAuth ecrit a la main, et plus sur. \\
Web service API XML \texttt{getZummHoneyActualQuantity} & API REST JSON (\texttt{@GetMapping}) + OpenAPI / Swagger & JSON au lieu de XML, documentation auto-generee ; XML reste possible via \texttt{produces = APPLICATION\_XML}. \\
X.509 / SSL (manuel) & Spring Boot SSL (\texttt{server.ssl.*}) + Keystore & Activation HTTPS declarative, memes certificats X.509, configuration triviale. \\
JavaScript (front libre) & React / Vue / Angular (ou Thymeleaf + JS) & Front decouple consommant l'API REST, parite web et mobile simplifiee. \\
\bottomrule
\end{longtable}

%------------------------------------------------------------------------------
\section{Tableau de correspondance : infrastructure et deploiement}
%------------------------------------------------------------------------------
\begin{longtable}{>{\bfseries}p{4cm} p{4.5cm} p{5.3cm}}
\toprule
\textbf{J2EE} & \textbf{Spring Boot} & \textbf{Benefice} \\
\midrule
\endhead
Serveur d'applications lourd (WildFly, GlassFish, WebLogic) & Tomcat / Jetty embarque & L'application se lance seule : \texttt{java -jar app.jar}, plus de serveur a installer. \\
Packaging WAR / EAR & JAR executable (fat-jar) & Un seul artefact autonome, ideal pour le cloud et Docker. \\
Deploiement manuel sur conteneur & Docker + \texttt{docker-compose} & Alignement DevOps du marche, integration continue facilitee. \\
SGBD relationnel via JDBC & PostgreSQL / MySQL via JPA + Flyway & Schema versionne, reproductible entre environnements. \\
Gestion des dependances manuelle (Ant) & Maven / Gradle + Spring Boot Starters & Une dependance = un starter (\texttt{-web}, \texttt{-security}, \texttt{-data-jpa}). \\
\bottomrule
\end{longtable}

%------------------------------------------------------------------------------
\section{Correspondance des couches}
%------------------------------------------------------------------------------
La vue en couches faiblement couplees exigee au chapitre~\ref{chap:conception}
(figure~\ref{fig:couches}) est \textbf{integralement preservee} : aucune couche
ne disparait, seule la mise en oeuvre est modernisee.

\begin{figure}[htbp]
\centering
\begin{tikzpicture}[
  j2ee/.style={draw=mielfonce,fill=grisdoux,thick,minimum width=5.4cm,minimum height=8mm,font=\footnotesize,align=center},
  spring/.style={draw=mielfonce,fill=miel!15,thick,minimum width=5.4cm,minimum height=8mm,font=\footnotesize,align=center},
  fl/.style={-{Latex[length=2mm]},mielfonce,thick}
]
% en-tetes
\node[font=\small\bfseries,mielfonce] at (0,5) {J2EE (impose)};
\node[font=\small\bfseries,mielfonce] at (7,5) {Spring Boot (cible)};
% couches J2EE
\node[j2ee] (p1) at (0,4) {Presentation : JSP + JavaScript};
\node[j2ee] (c1) at (0,3) {Controle : Servlets (MVC)};
\node[j2ee] (m1) at (0,2) {Metier : EJB};
\node[j2ee] (d1) at (0,1) {Acces donnees : JDBC};
\node[j2ee] (b1) at (0,0) {Base de donnees relationnelle};
% couches Spring
\node[spring] (p2) at (7,4) {@Controller + Thymeleaf / React};
\node[spring] (c2) at (7,3) {@RestController (Spring MVC)};
\node[spring] (m2) at (7,2) {@Service (Spring Beans)};
\node[spring] (d2) at (7,1) {@Repository (Spring Data JPA)};
\node[spring] (b2) at (7,0) {PostgreSQL / MySQL (inchange)};
% fleches de correspondance
\draw[fl] (p1)--(p2); \draw[fl] (c1)--(c2); \draw[fl] (m1)--(m2);
\draw[fl] (d1)--(d2); \draw[fl] (b1)--(b2);
\end{tikzpicture}
\caption{Correspondance des couches J2EE vers Spring Boot : architecture preservee, mise en oeuvre modernisee.}
\label{fig:migration-couches}
\end{figure}

Les modules transverses suivent la meme logique : i18n XML \textrightarrow{}
\texttt{MessageSource}, \texttt{ConfigZumm.ini} \textrightarrow{}
\texttt{application.yml}, service web XML \textrightarrow{} API REST + OpenAPI.

%------------------------------------------------------------------------------
\section{Argumentaire : conformite academique et vision marche}
%------------------------------------------------------------------------------
\textbf{Z\"umm} repose, conformement au cahier des charges, sur une
architecture \textbf{J2EE multi-niveaux, evolutive et faiblement couplee} : MVC,
Servlets, JSP, EJB, JDBC, internationalisation par fichier externalise, services
web securises (OAuth, X.509 / SSL) et API interrogeable par des applications
tierces. Ce socle a ete retenu et implemente \textbf{tel qu'exige}, car il
constitue le referentiel d'evaluation du projet et materialise les competences
fondamentales du parcours DSI : maitrise des couches d'un systeme d'information,
separation des responsabilites et interoperabilite.

Une analyse \textbf{orientee marche} conduit toutefois a un constat important.
Les briques Java EE historiques (JSP, Servlets \og nus \fg{} et EJB) sont
aujourd'hui considerees comme \emph{heritees}. Les recrutements sur les nouveaux
systemes d'information privilegient tres majoritairement l'ecosysteme
\textbf{Spring Boot}, devenu le standard \emph{de facto} du developpement Java
d'entreprise. Ignorer cette realite reviendrait a livrer une solution
techniquement conforme mais \textbf{decalee des pratiques professionnelles
actuelles}.

Nous avons donc adopte une \textbf{demarche a deux niveaux de lecture}, sans
jamais sortir du langage \textbf{Java}, ce qui garantit la continuite des
competences acquises :
\begin{enumerate}
  \item \textbf{Niveau conformite} : le livrable respecte integralement les
        technologies imposees et la grille d'evaluation. Les diagrammes de
        conception (sequence Agent \textrightarrow{} JSP \textrightarrow{}
        Servlet \textrightarrow{} EJB \textrightarrow{} JDBC, couches,
        deploiement) refletent fidelement cette architecture.
  \item \textbf{Niveau modernisation} : pour chaque technologie imposee, son
        equivalent Spring Boot est demontre (Servlet \textrightarrow{}
        \texttt{@RestController}, EJB \textrightarrow{} \texttt{@Service}, JDBC
        \textrightarrow{} Spring Data JPA, XML i18n \textrightarrow{}
        \texttt{MessageSource}, \texttt{ConfigZumm.ini} \textrightarrow{}
        \texttt{application.yml}, WS XML \textrightarrow{} API REST + OpenAPI,
        X.509 / SSL \textrightarrow{} configuration SSL declarative). Cette
        correspondance prouve que \textbf{l'architecture cible reste identique} :
        seule la mise en oeuvre gagne en concision, en securite et en
        maintenabilite.
\end{enumerate}

Ce choix presente un troisieme avantage, en phase directe avec les debouches du
parcours DSI. Le cahier des charges prevoit un \textbf{volet capteurs et analyse
predictive} (detection d'essaimage, traitement acoustique). Ce module se prete
naturellement a un \textbf{microservice Python} (pandas, scikit-learn) expose au
back-end Java via REST, illustrant concretement le metier d'\textbf{analyste de
donnees} cite au plan d'etudes et l'ouverture d'une \textbf{architecture
faiblement couplee} a l'heterogeneite technologique.

\begin{encadre}
\textbf{Conclusion :} le projet satisfait pleinement les exigences academiques
(J2EE) tout en documentant sa \textbf{trajectoire de modernisation} (Spring Boot
et microservice data). Cette double lecture transforme une contrainte
pedagogique en atout professionnel : elle demontre non seulement la maitrise des
fondamentaux du systeme d'information, mais aussi la capacite a faire evoluer une
architecture vers les standards recherches sur le marche de l'emploi en 2026.
\end{encadre} dans cahier_des_charges.tex
%  (utilise la charte "Miel & noir" : miel, mielfonce, encadre, colonne C,
%   longtable, booktabs -- deja definis dans le preambule principal)
%==============================================================================

%==============================================================================
\chapter{Annexe - Choix technologique : academique et vision marche}
\label{chap:migration}
%==============================================================================
Cette annexe repond a la question du \emph{pourquoi} et du \emph{comment} du
langage de developpement. Elle confronte le socle \textbf{J2EE} impose par le
present cahier des charges a l'ecosysteme \textbf{Spring Boot}, devenu le
standard du developpement Java d'entreprise, et documente une trajectoire de
modernisation sans changement de langage.

\section{Positionnement}
Le systeme \textbf{Z\"umm} est implemente en \textbf{J2EE (Java EE)}
conformement aux exigences techniques (chapitre~\ref{chap:conception} et grille
d'evaluation) : MVC, Servlets, JSP, EJB, JDBC, internationalisation
externalisee, services web securises (OAuth, X.509 / SSL) et API interrogeable
par des applications tierces. Ce socle materialise les competences fondamentales
du parcours DSI : maitrise des couches d'un systeme d'information, separation des
responsabilites et interoperabilite.

\begin{encadre}
\textbf{Constat marche :} les briques Java EE historiques (JSP, Servlets
\og nus \fg{} et EJB) sont aujourd'hui considerees comme \emph{heritees}
(\emph{legacy}). Les recrutements sur les nouveaux systemes d'information
privilegient tres majoritairement l'ecosysteme \textbf{Spring Boot}. La presente
annexe conserve le langage \textbf{Java} et demontre, techno par techno,
l'equivalent moderne de chaque exigence.
\end{encadre}

%------------------------------------------------------------------------------
\section{Tableau de correspondance : coeur technique}
%------------------------------------------------------------------------------
Le tableau suivant met en regard chaque technologie imposee (et notee) avec son
equivalent Spring Boot, ainsi que la nature concrete du changement.

\begin{longtable}{>{\bfseries}p{3cm} p{3.4cm} p{7.2cm}}
\toprule
\textbf{J2EE (impose)} & \textbf{Spring Boot} & \textbf{Ce qui change concretement} \\
\midrule
\endhead
MVC (manuel) & Spring MVC (\texttt{@Controller}, \texttt{@RequestMapping}) & Le pattern reste ; le routage devient declaratif par annotations au lieu du \texttt{web.xml}. \\
Servlet (\texttt{HttpServlet}, \texttt{doGet/doPost}) & \texttt{@RestController} / \texttt{@Controller} & Plus de \texttt{extends HttpServlet} : une methode annotee = un endpoint. Le \texttt{DispatcherServlet} orchestre en arriere-plan. \\
JSP (presentation) & Thymeleaf (server-side) ou React / Angular (SPA) & Templating moderne sans scriptlets Java ; pour un effet marche fort, API REST + front decouple. \\
EJB (\texttt{@Stateless}, conteneur EJB) & Spring \texttt{@Service} / \texttt{@Component} & Meme role (logique metier, injection, transactions) sans conteneur EJB lourd ; \texttt{@Transactional} gere les transactions. \\
JDBC (SQL manuel, \texttt{ResultSet}) & Spring Data JPA / Hibernate & Le mapping objet-relationnel remplace le SQL a la main ; \texttt{JpaRepository} genere les CRUD (\texttt{JdbcTemplate} reste disponible). \\
i18n par fichier XML & \texttt{messages\_fr/en/ar.properties} + \texttt{MessageSource} & Meme principe d'externalisation des textes, format \texttt{.properties} standard, toujours FR / EN / AR. \\
\texttt{ConfigZumm.ini} & \texttt{application.yml} + \texttt{@ConfigurationProperties} & Parametrage sans recompilation, injectable par profil ; les seuils et unites deviennent des beans types. \\
OAuth Google (WS XML maison) & Spring Security OAuth2 Client & Login Google en quelques lignes de configuration au lieu d'un client OAuth ecrit a la main, et plus sur. \\
Web service API XML \texttt{getZummHoneyActualQuantity} & API REST JSON (\texttt{@GetMapping}) + OpenAPI / Swagger & JSON au lieu de XML, documentation auto-generee ; XML reste possible via \texttt{produces = APPLICATION\_XML}. \\
X.509 / SSL (manuel) & Spring Boot SSL (\texttt{server.ssl.*}) + Keystore & Activation HTTPS declarative, memes certificats X.509, configuration triviale. \\
JavaScript (front libre) & React / Vue / Angular (ou Thymeleaf + JS) & Front decouple consommant l'API REST, parite web et mobile simplifiee. \\
\bottomrule
\end{longtable}

%------------------------------------------------------------------------------
\section{Tableau de correspondance : infrastructure et deploiement}
%------------------------------------------------------------------------------
\begin{longtable}{>{\bfseries}p{4cm} p{4.5cm} p{5.3cm}}
\toprule
\textbf{J2EE} & \textbf{Spring Boot} & \textbf{Benefice} \\
\midrule
\endhead
Serveur d'applications lourd (WildFly, GlassFish, WebLogic) & Tomcat / Jetty embarque & L'application se lance seule : \texttt{java -jar app.jar}, plus de serveur a installer. \\
Packaging WAR / EAR & JAR executable (fat-jar) & Un seul artefact autonome, ideal pour le cloud et Docker. \\
Deploiement manuel sur conteneur & Docker + \texttt{docker-compose} & Alignement DevOps du marche, integration continue facilitee. \\
SGBD relationnel via JDBC & PostgreSQL / MySQL via JPA + Flyway & Schema versionne, reproductible entre environnements. \\
Gestion des dependances manuelle (Ant) & Maven / Gradle + Spring Boot Starters & Une dependance = un starter (\texttt{-web}, \texttt{-security}, \texttt{-data-jpa}). \\
\bottomrule
\end{longtable}

%------------------------------------------------------------------------------
\section{Correspondance des couches}
%------------------------------------------------------------------------------
La vue en couches faiblement couplees exigee au chapitre~\ref{chap:conception}
(figure~\ref{fig:couches}) est \textbf{integralement preservee} : aucune couche
ne disparait, seule la mise en oeuvre est modernisee.

\begin{figure}[htbp]
\centering
\begin{tikzpicture}[
  j2ee/.style={draw=mielfonce,fill=grisdoux,thick,minimum width=5.4cm,minimum height=8mm,font=\footnotesize,align=center},
  spring/.style={draw=mielfonce,fill=miel!15,thick,minimum width=5.4cm,minimum height=8mm,font=\footnotesize,align=center},
  fl/.style={-{Latex[length=2mm]},mielfonce,thick}
]
% en-tetes
\node[font=\small\bfseries,mielfonce] at (0,5) {J2EE (impose)};
\node[font=\small\bfseries,mielfonce] at (7,5) {Spring Boot (cible)};
% couches J2EE
\node[j2ee] (p1) at (0,4) {Presentation : JSP + JavaScript};
\node[j2ee] (c1) at (0,3) {Controle : Servlets (MVC)};
\node[j2ee] (m1) at (0,2) {Metier : EJB};
\node[j2ee] (d1) at (0,1) {Acces donnees : JDBC};
\node[j2ee] (b1) at (0,0) {Base de donnees relationnelle};
% couches Spring
\node[spring] (p2) at (7,4) {@Controller + Thymeleaf / React};
\node[spring] (c2) at (7,3) {@RestController (Spring MVC)};
\node[spring] (m2) at (7,2) {@Service (Spring Beans)};
\node[spring] (d2) at (7,1) {@Repository (Spring Data JPA)};
\node[spring] (b2) at (7,0) {PostgreSQL / MySQL (inchange)};
% fleches de correspondance
\draw[fl] (p1)--(p2); \draw[fl] (c1)--(c2); \draw[fl] (m1)--(m2);
\draw[fl] (d1)--(d2); \draw[fl] (b1)--(b2);
\end{tikzpicture}
\caption{Correspondance des couches J2EE vers Spring Boot : architecture preservee, mise en oeuvre modernisee.}
\label{fig:migration-couches}
\end{figure}

Les modules transverses suivent la meme logique : i18n XML \textrightarrow{}
\texttt{MessageSource}, \texttt{ConfigZumm.ini} \textrightarrow{}
\texttt{application.yml}, service web XML \textrightarrow{} API REST + OpenAPI.

%------------------------------------------------------------------------------
\section{Argumentaire : conformite academique et vision marche}
%------------------------------------------------------------------------------
\textbf{Z\"umm} repose, conformement au cahier des charges, sur une
architecture \textbf{J2EE multi-niveaux, evolutive et faiblement couplee} : MVC,
Servlets, JSP, EJB, JDBC, internationalisation par fichier externalise, services
web securises (OAuth, X.509 / SSL) et API interrogeable par des applications
tierces. Ce socle a ete retenu et implemente \textbf{tel qu'exige}, car il
constitue le referentiel d'evaluation du projet et materialise les competences
fondamentales du parcours DSI : maitrise des couches d'un systeme d'information,
separation des responsabilites et interoperabilite.

Une analyse \textbf{orientee marche} conduit toutefois a un constat important.
Les briques Java EE historiques (JSP, Servlets \og nus \fg{} et EJB) sont
aujourd'hui considerees comme \emph{heritees}. Les recrutements sur les nouveaux
systemes d'information privilegient tres majoritairement l'ecosysteme
\textbf{Spring Boot}, devenu le standard \emph{de facto} du developpement Java
d'entreprise. Ignorer cette realite reviendrait a livrer une solution
techniquement conforme mais \textbf{decalee des pratiques professionnelles
actuelles}.

Nous avons donc adopte une \textbf{demarche a deux niveaux de lecture}, sans
jamais sortir du langage \textbf{Java}, ce qui garantit la continuite des
competences acquises :
\begin{enumerate}
  \item \textbf{Niveau conformite} : le livrable respecte integralement les
        technologies imposees et la grille d'evaluation. Les diagrammes de
        conception (sequence Agent \textrightarrow{} JSP \textrightarrow{}
        Servlet \textrightarrow{} EJB \textrightarrow{} JDBC, couches,
        deploiement) refletent fidelement cette architecture.
  \item \textbf{Niveau modernisation} : pour chaque technologie imposee, son
        equivalent Spring Boot est demontre (Servlet \textrightarrow{}
        \texttt{@RestController}, EJB \textrightarrow{} \texttt{@Service}, JDBC
        \textrightarrow{} Spring Data JPA, XML i18n \textrightarrow{}
        \texttt{MessageSource}, \texttt{ConfigZumm.ini} \textrightarrow{}
        \texttt{application.yml}, WS XML \textrightarrow{} API REST + OpenAPI,
        X.509 / SSL \textrightarrow{} configuration SSL declarative). Cette
        correspondance prouve que \textbf{l'architecture cible reste identique} :
        seule la mise en oeuvre gagne en concision, en securite et en
        maintenabilite.
\end{enumerate}

Ce choix presente un troisieme avantage, en phase directe avec les debouches du
parcours DSI. Le cahier des charges prevoit un \textbf{volet capteurs et analyse
predictive} (detection d'essaimage, traitement acoustique). Ce module se prete
naturellement a un \textbf{microservice Python} (pandas, scikit-learn) expose au
back-end Java via REST, illustrant concretement le metier d'\textbf{analyste de
donnees} cite au plan d'etudes et l'ouverture d'une \textbf{architecture
faiblement couplee} a l'heterogeneite technologique.

\begin{encadre}
\textbf{Conclusion :} le projet satisfait pleinement les exigences academiques
(J2EE) tout en documentant sa \textbf{trajectoire de modernisation} (Spring Boot
et microservice data). Cette double lecture transforme une contrainte
pedagogique en atout professionnel : elle demontre non seulement la maitrise des
fondamentaux du systeme d'information, mais aussi la capacite a faire evoluer une
architecture vers les standards recherches sur le marche de l'emploi en 2026.
\end{encadre} dans cahier_des_charges.tex
%  (utilise la charte "Miel & noir" : miel, mielfonce, encadre, colonne C,
%   longtable, booktabs -- deja definis dans le preambule principal)
%==============================================================================

%==============================================================================
\chapter{Annexe - Choix technologique : academique et vision marche}
\label{chap:migration}
%==============================================================================
Cette annexe repond a la question du \emph{pourquoi} et du \emph{comment} du
langage de developpement. Elle confronte le socle \textbf{J2EE} impose par le
present cahier des charges a l'ecosysteme \textbf{Spring Boot}, devenu le
standard du developpement Java d'entreprise, et documente une trajectoire de
modernisation sans changement de langage.

\section{Positionnement}
Le systeme \textbf{Z\"umm} est implemente en \textbf{J2EE (Java EE)}
conformement aux exigences techniques (chapitre~\ref{chap:conception} et grille
d'evaluation) : MVC, Servlets, JSP, EJB, JDBC, internationalisation
externalisee, services web securises (OAuth, X.509 / SSL) et API interrogeable
par des applications tierces. Ce socle materialise les competences fondamentales
du parcours DSI : maitrise des couches d'un systeme d'information, separation des
responsabilites et interoperabilite.

\begin{encadre}
\textbf{Constat marche :} les briques Java EE historiques (JSP, Servlets
\og nus \fg{} et EJB) sont aujourd'hui considerees comme \emph{heritees}
(\emph{legacy}). Les recrutements sur les nouveaux systemes d'information
privilegient tres majoritairement l'ecosysteme \textbf{Spring Boot}. La presente
annexe conserve le langage \textbf{Java} et demontre, techno par techno,
l'equivalent moderne de chaque exigence.
\end{encadre}

%------------------------------------------------------------------------------
\section{Tableau de correspondance : coeur technique}
%------------------------------------------------------------------------------
Le tableau suivant met en regard chaque technologie imposee (et notee) avec son
equivalent Spring Boot, ainsi que la nature concrete du changement.

\begin{longtable}{>{\bfseries}p{3cm} p{3.4cm} p{7.2cm}}
\toprule
\textbf{J2EE (impose)} & \textbf{Spring Boot} & \textbf{Ce qui change concretement} \\
\midrule
\endhead
MVC (manuel) & Spring MVC (\texttt{@Controller}, \texttt{@RequestMapping}) & Le pattern reste ; le routage devient declaratif par annotations au lieu du \texttt{web.xml}. \\
Servlet (\texttt{HttpServlet}, \texttt{doGet/doPost}) & \texttt{@RestController} / \texttt{@Controller} & Plus de \texttt{extends HttpServlet} : une methode annotee = un endpoint. Le \texttt{DispatcherServlet} orchestre en arriere-plan. \\
JSP (presentation) & Thymeleaf (server-side) ou React / Angular (SPA) & Templating moderne sans scriptlets Java ; pour un effet marche fort, API REST + front decouple. \\
EJB (\texttt{@Stateless}, conteneur EJB) & Spring \texttt{@Service} / \texttt{@Component} & Meme role (logique metier, injection, transactions) sans conteneur EJB lourd ; \texttt{@Transactional} gere les transactions. \\
JDBC (SQL manuel, \texttt{ResultSet}) & Spring Data JPA / Hibernate & Le mapping objet-relationnel remplace le SQL a la main ; \texttt{JpaRepository} genere les CRUD (\texttt{JdbcTemplate} reste disponible). \\
i18n par fichier XML & \texttt{messages\_fr/en/ar.properties} + \texttt{MessageSource} & Meme principe d'externalisation des textes, format \texttt{.properties} standard, toujours FR / EN / AR. \\
\texttt{ConfigZumm.ini} & \texttt{application.yml} + \texttt{@ConfigurationProperties} & Parametrage sans recompilation, injectable par profil ; les seuils et unites deviennent des beans types. \\
OAuth Google (WS XML maison) & Spring Security OAuth2 Client & Login Google en quelques lignes de configuration au lieu d'un client OAuth ecrit a la main, et plus sur. \\
Web service API XML \texttt{getZummHoneyActualQuantity} & API REST JSON (\texttt{@GetMapping}) + OpenAPI / Swagger & JSON au lieu de XML, documentation auto-generee ; XML reste possible via \texttt{produces = APPLICATION\_XML}. \\
X.509 / SSL (manuel) & Spring Boot SSL (\texttt{server.ssl.*}) + Keystore & Activation HTTPS declarative, memes certificats X.509, configuration triviale. \\
JavaScript (front libre) & React / Vue / Angular (ou Thymeleaf + JS) & Front decouple consommant l'API REST, parite web et mobile simplifiee. \\
\bottomrule
\end{longtable}

%------------------------------------------------------------------------------
\section{Tableau de correspondance : infrastructure et deploiement}
%------------------------------------------------------------------------------
\begin{longtable}{>{\bfseries}p{4cm} p{4.5cm} p{5.3cm}}
\toprule
\textbf{J2EE} & \textbf{Spring Boot} & \textbf{Benefice} \\
\midrule
\endhead
Serveur d'applications lourd (WildFly, GlassFish, WebLogic) & Tomcat / Jetty embarque & L'application se lance seule : \texttt{java -jar app.jar}, plus de serveur a installer. \\
Packaging WAR / EAR & JAR executable (fat-jar) & Un seul artefact autonome, ideal pour le cloud et Docker. \\
Deploiement manuel sur conteneur & Docker + \texttt{docker-compose} & Alignement DevOps du marche, integration continue facilitee. \\
SGBD relationnel via JDBC & PostgreSQL / MySQL via JPA + Flyway & Schema versionne, reproductible entre environnements. \\
Gestion des dependances manuelle (Ant) & Maven / Gradle + Spring Boot Starters & Une dependance = un starter (\texttt{-web}, \texttt{-security}, \texttt{-data-jpa}). \\
\bottomrule
\end{longtable}

%------------------------------------------------------------------------------
\section{Correspondance des couches}
%------------------------------------------------------------------------------
La vue en couches faiblement couplees exigee au chapitre~\ref{chap:conception}
(figure~\ref{fig:couches}) est \textbf{integralement preservee} : aucune couche
ne disparait, seule la mise en oeuvre est modernisee.

\begin{figure}[htbp]
\centering
\begin{tikzpicture}[
  j2ee/.style={draw=mielfonce,fill=grisdoux,thick,minimum width=5.4cm,minimum height=8mm,font=\footnotesize,align=center},
  spring/.style={draw=mielfonce,fill=miel!15,thick,minimum width=5.4cm,minimum height=8mm,font=\footnotesize,align=center},
  fl/.style={-{Latex[length=2mm]},mielfonce,thick}
]
% en-tetes
\node[font=\small\bfseries,mielfonce] at (0,5) {J2EE (impose)};
\node[font=\small\bfseries,mielfonce] at (7,5) {Spring Boot (cible)};
% couches J2EE
\node[j2ee] (p1) at (0,4) {Presentation : JSP + JavaScript};
\node[j2ee] (c1) at (0,3) {Controle : Servlets (MVC)};
\node[j2ee] (m1) at (0,2) {Metier : EJB};
\node[j2ee] (d1) at (0,1) {Acces donnees : JDBC};
\node[j2ee] (b1) at (0,0) {Base de donnees relationnelle};
% couches Spring
\node[spring] (p2) at (7,4) {@Controller + Thymeleaf / React};
\node[spring] (c2) at (7,3) {@RestController (Spring MVC)};
\node[spring] (m2) at (7,2) {@Service (Spring Beans)};
\node[spring] (d2) at (7,1) {@Repository (Spring Data JPA)};
\node[spring] (b2) at (7,0) {PostgreSQL / MySQL (inchange)};
% fleches de correspondance
\draw[fl] (p1)--(p2); \draw[fl] (c1)--(c2); \draw[fl] (m1)--(m2);
\draw[fl] (d1)--(d2); \draw[fl] (b1)--(b2);
\end{tikzpicture}
\caption{Correspondance des couches J2EE vers Spring Boot : architecture preservee, mise en oeuvre modernisee.}
\label{fig:migration-couches}
\end{figure}

Les modules transverses suivent la meme logique : i18n XML \textrightarrow{}
\texttt{MessageSource}, \texttt{ConfigZumm.ini} \textrightarrow{}
\texttt{application.yml}, service web XML \textrightarrow{} API REST + OpenAPI.

%------------------------------------------------------------------------------
\section{Argumentaire : conformite academique et vision marche}
%------------------------------------------------------------------------------
\textbf{Z\"umm} repose, conformement au cahier des charges, sur une
architecture \textbf{J2EE multi-niveaux, evolutive et faiblement couplee} : MVC,
Servlets, JSP, EJB, JDBC, internationalisation par fichier externalise, services
web securises (OAuth, X.509 / SSL) et API interrogeable par des applications
tierces. Ce socle a ete retenu et implemente \textbf{tel qu'exige}, car il
constitue le referentiel d'evaluation du projet et materialise les competences
fondamentales du parcours DSI : maitrise des couches d'un systeme d'information,
separation des responsabilites et interoperabilite.

Une analyse \textbf{orientee marche} conduit toutefois a un constat important.
Les briques Java EE historiques (JSP, Servlets \og nus \fg{} et EJB) sont
aujourd'hui considerees comme \emph{heritees}. Les recrutements sur les nouveaux
systemes d'information privilegient tres majoritairement l'ecosysteme
\textbf{Spring Boot}, devenu le standard \emph{de facto} du developpement Java
d'entreprise. Ignorer cette realite reviendrait a livrer une solution
techniquement conforme mais \textbf{decalee des pratiques professionnelles
actuelles}.

Nous avons donc adopte une \textbf{demarche a deux niveaux de lecture}, sans
jamais sortir du langage \textbf{Java}, ce qui garantit la continuite des
competences acquises :
\begin{enumerate}
  \item \textbf{Niveau conformite} : le livrable respecte integralement les
        technologies imposees et la grille d'evaluation. Les diagrammes de
        conception (sequence Agent \textrightarrow{} JSP \textrightarrow{}
        Servlet \textrightarrow{} EJB \textrightarrow{} JDBC, couches,
        deploiement) refletent fidelement cette architecture.
  \item \textbf{Niveau modernisation} : pour chaque technologie imposee, son
        equivalent Spring Boot est demontre (Servlet \textrightarrow{}
        \texttt{@RestController}, EJB \textrightarrow{} \texttt{@Service}, JDBC
        \textrightarrow{} Spring Data JPA, XML i18n \textrightarrow{}
        \texttt{MessageSource}, \texttt{ConfigZumm.ini} \textrightarrow{}
        \texttt{application.yml}, WS XML \textrightarrow{} API REST + OpenAPI,
        X.509 / SSL \textrightarrow{} configuration SSL declarative). Cette
        correspondance prouve que \textbf{l'architecture cible reste identique} :
        seule la mise en oeuvre gagne en concision, en securite et en
        maintenabilite.
\end{enumerate}

Ce choix presente un troisieme avantage, en phase directe avec les debouches du
parcours DSI. Le cahier des charges prevoit un \textbf{volet capteurs et analyse
predictive} (detection d'essaimage, traitement acoustique). Ce module se prete
naturellement a un \textbf{microservice Python} (pandas, scikit-learn) expose au
back-end Java via REST, illustrant concretement le metier d'\textbf{analyste de
donnees} cite au plan d'etudes et l'ouverture d'une \textbf{architecture
faiblement couplee} a l'heterogeneite technologique.

\begin{encadre}
\textbf{Conclusion :} le projet satisfait pleinement les exigences academiques
(J2EE) tout en documentant sa \textbf{trajectoire de modernisation} (Spring Boot
et microservice data). Cette double lecture transforme une contrainte
pedagogique en atout professionnel : elle demontre non seulement la maitrise des
fondamentaux du systeme d'information, mais aussi la capacite a faire evoluer une
architecture vers les standards recherches sur le marche de l'emploi en 2026.
\end{encadre} dans cahier_des_charges.tex
%  (utilise la charte "Miel & noir" : miel, mielfonce, encadre, colonne C,
%   longtable, booktabs -- deja definis dans le preambule principal)
%==============================================================================

%==============================================================================
\chapter{Annexe - Choix technologique : academique et vision marche}
\label{chap:migration}
%==============================================================================
Cette annexe repond a la question du \emph{pourquoi} et du \emph{comment} du
langage de developpement. Elle confronte le socle \textbf{J2EE} impose par le
present cahier des charges a l'ecosysteme \textbf{Spring Boot}, devenu le
standard du developpement Java d'entreprise, et documente une trajectoire de
modernisation sans changement de langage.

\section{Positionnement}
Le systeme \textbf{Z\"umm} est implemente en \textbf{J2EE (Java EE)}
conformement aux exigences techniques (chapitre~\ref{chap:conception} et grille
d'evaluation) : MVC, Servlets, JSP, EJB, JDBC, internationalisation
externalisee, services web securises (OAuth, X.509 / SSL) et API interrogeable
par des applications tierces. Ce socle materialise les competences fondamentales
du parcours DSI : maitrise des couches d'un systeme d'information, separation des
responsabilites et interoperabilite.

\begin{encadre}
\textbf{Constat marche :} les briques Java EE historiques (JSP, Servlets
\og nus \fg{} et EJB) sont aujourd'hui considerees comme \emph{heritees}
(\emph{legacy}). Les recrutements sur les nouveaux systemes d'information
privilegient tres majoritairement l'ecosysteme \textbf{Spring Boot}. La presente
annexe conserve le langage \textbf{Java} et demontre, techno par techno,
l'equivalent moderne de chaque exigence.
\end{encadre}

%------------------------------------------------------------------------------
\section{Tableau de correspondance : coeur technique}
%------------------------------------------------------------------------------
Le tableau suivant met en regard chaque technologie imposee (et notee) avec son
equivalent Spring Boot, ainsi que la nature concrete du changement.

\begin{longtable}{>{\bfseries}p{3cm} p{3.4cm} p{7.2cm}}
\toprule
\textbf{J2EE (impose)} & \textbf{Spring Boot} & \textbf{Ce qui change concretement} \\
\midrule
\endhead
MVC (manuel) & Spring MVC (\texttt{@Controller}, \texttt{@RequestMapping}) & Le pattern reste ; le routage devient declaratif par annotations au lieu du \texttt{web.xml}. \\
Servlet (\texttt{HttpServlet}, \texttt{doGet/doPost}) & \texttt{@RestController} / \texttt{@Controller} & Plus de \texttt{extends HttpServlet} : une methode annotee = un endpoint. Le \texttt{DispatcherServlet} orchestre en arriere-plan. \\
JSP (presentation) & Thymeleaf (server-side) ou React / Angular (SPA) & Templating moderne sans scriptlets Java ; pour un effet marche fort, API REST + front decouple. \\
EJB (\texttt{@Stateless}, conteneur EJB) & Spring \texttt{@Service} / \texttt{@Component} & Meme role (logique metier, injection, transactions) sans conteneur EJB lourd ; \texttt{@Transactional} gere les transactions. \\
JDBC (SQL manuel, \texttt{ResultSet}) & Spring Data JPA / Hibernate & Le mapping objet-relationnel remplace le SQL a la main ; \texttt{JpaRepository} genere les CRUD (\texttt{JdbcTemplate} reste disponible). \\
i18n par fichier XML & \texttt{messages\_fr/en/ar.properties} + \texttt{MessageSource} & Meme principe d'externalisation des textes, format \texttt{.properties} standard, toujours FR / EN / AR. \\
\texttt{ConfigZumm.ini} & \texttt{application.yml} + \texttt{@ConfigurationProperties} & Parametrage sans recompilation, injectable par profil ; les seuils et unites deviennent des beans types. \\
OAuth Google (WS XML maison) & Spring Security OAuth2 Client & Login Google en quelques lignes de configuration au lieu d'un client OAuth ecrit a la main, et plus sur. \\
Web service API XML \texttt{getZummHoneyActualQuantity} & API REST JSON (\texttt{@GetMapping}) + OpenAPI / Swagger & JSON au lieu de XML, documentation auto-generee ; XML reste possible via \texttt{produces = APPLICATION\_XML}. \\
X.509 / SSL (manuel) & Spring Boot SSL (\texttt{server.ssl.*}) + Keystore & Activation HTTPS declarative, memes certificats X.509, configuration triviale. \\
JavaScript (front libre) & React / Vue / Angular (ou Thymeleaf + JS) & Front decouple consommant l'API REST, parite web et mobile simplifiee. \\
\bottomrule
\end{longtable}

%------------------------------------------------------------------------------
\section{Tableau de correspondance : infrastructure et deploiement}
%------------------------------------------------------------------------------
\begin{longtable}{>{\bfseries}p{4cm} p{4.5cm} p{5.3cm}}
\toprule
\textbf{J2EE} & \textbf{Spring Boot} & \textbf{Benefice} \\
\midrule
\endhead
Serveur d'applications lourd (WildFly, GlassFish, WebLogic) & Tomcat / Jetty embarque & L'application se lance seule : \texttt{java -jar app.jar}, plus de serveur a installer. \\
Packaging WAR / EAR & JAR executable (fat-jar) & Un seul artefact autonome, ideal pour le cloud et Docker. \\
Deploiement manuel sur conteneur & Docker + \texttt{docker-compose} & Alignement DevOps du marche, integration continue facilitee. \\
SGBD relationnel via JDBC & PostgreSQL / MySQL via JPA + Flyway & Schema versionne, reproductible entre environnements. \\
Gestion des dependances manuelle (Ant) & Maven / Gradle + Spring Boot Starters & Une dependance = un starter (\texttt{-web}, \texttt{-security}, \texttt{-data-jpa}). \\
\bottomrule
\end{longtable}

%------------------------------------------------------------------------------
\section{Correspondance des couches}
%------------------------------------------------------------------------------
La vue en couches faiblement couplees exigee au chapitre~\ref{chap:conception}
(figure~\ref{fig:couches}) est \textbf{integralement preservee} : aucune couche
ne disparait, seule la mise en oeuvre est modernisee.

\begin{figure}[htbp]
\centering
\begin{tikzpicture}[
  j2ee/.style={draw=mielfonce,fill=grisdoux,thick,minimum width=5.4cm,minimum height=8mm,font=\footnotesize,align=center},
  spring/.style={draw=mielfonce,fill=miel!15,thick,minimum width=5.4cm,minimum height=8mm,font=\footnotesize,align=center},
  fl/.style={-{Latex[length=2mm]},mielfonce,thick}
]
% en-tetes
\node[font=\small\bfseries,mielfonce] at (0,5) {J2EE (impose)};
\node[font=\small\bfseries,mielfonce] at (7,5) {Spring Boot (cible)};
% couches J2EE
\node[j2ee] (p1) at (0,4) {Presentation : JSP + JavaScript};
\node[j2ee] (c1) at (0,3) {Controle : Servlets (MVC)};
\node[j2ee] (m1) at (0,2) {Metier : EJB};
\node[j2ee] (d1) at (0,1) {Acces donnees : JDBC};
\node[j2ee] (b1) at (0,0) {Base de donnees relationnelle};
% couches Spring
\node[spring] (p2) at (7,4) {@Controller + Thymeleaf / React};
\node[spring] (c2) at (7,3) {@RestController (Spring MVC)};
\node[spring] (m2) at (7,2) {@Service (Spring Beans)};
\node[spring] (d2) at (7,1) {@Repository (Spring Data JPA)};
\node[spring] (b2) at (7,0) {PostgreSQL / MySQL (inchange)};
% fleches de correspondance
\draw[fl] (p1)--(p2); \draw[fl] (c1)--(c2); \draw[fl] (m1)--(m2);
\draw[fl] (d1)--(d2); \draw[fl] (b1)--(b2);
\end{tikzpicture}
\caption{Correspondance des couches J2EE vers Spring Boot : architecture preservee, mise en oeuvre modernisee.}
\label{fig:migration-couches}
\end{figure}

Les modules transverses suivent la meme logique : i18n XML \textrightarrow{}
\texttt{MessageSource}, \texttt{ConfigZumm.ini} \textrightarrow{}
\texttt{application.yml}, service web XML \textrightarrow{} API REST + OpenAPI.

%------------------------------------------------------------------------------
\section{Argumentaire : conformite academique et vision marche}
%------------------------------------------------------------------------------
\textbf{Z\"umm} repose, conformement au cahier des charges, sur une
architecture \textbf{J2EE multi-niveaux, evolutive et faiblement couplee} : MVC,
Servlets, JSP, EJB, JDBC, internationalisation par fichier externalise, services
web securises (OAuth, X.509 / SSL) et API interrogeable par des applications
tierces. Ce socle a ete retenu et implemente \textbf{tel qu'exige}, car il
constitue le referentiel d'evaluation du projet et materialise les competences
fondamentales du parcours DSI : maitrise des couches d'un systeme d'information,
separation des responsabilites et interoperabilite.

Une analyse \textbf{orientee marche} conduit toutefois a un constat important.
Les briques Java EE historiques (JSP, Servlets \og nus \fg{} et EJB) sont
aujourd'hui considerees comme \emph{heritees}. Les recrutements sur les nouveaux
systemes d'information privilegient tres majoritairement l'ecosysteme
\textbf{Spring Boot}, devenu le standard \emph{de facto} du developpement Java
d'entreprise. Ignorer cette realite reviendrait a livrer une solution
techniquement conforme mais \textbf{decalee des pratiques professionnelles
actuelles}.

Nous avons donc adopte une \textbf{demarche a deux niveaux de lecture}, sans
jamais sortir du langage \textbf{Java}, ce qui garantit la continuite des
competences acquises :
\begin{enumerate}
  \item \textbf{Niveau conformite} : le livrable respecte integralement les
        technologies imposees et la grille d'evaluation. Les diagrammes de
        conception (sequence Agent \textrightarrow{} JSP \textrightarrow{}
        Servlet \textrightarrow{} EJB \textrightarrow{} JDBC, couches,
        deploiement) refletent fidelement cette architecture.
  \item \textbf{Niveau modernisation} : pour chaque technologie imposee, son
        equivalent Spring Boot est demontre (Servlet \textrightarrow{}
        \texttt{@RestController}, EJB \textrightarrow{} \texttt{@Service}, JDBC
        \textrightarrow{} Spring Data JPA, XML i18n \textrightarrow{}
        \texttt{MessageSource}, \texttt{ConfigZumm.ini} \textrightarrow{}
        \texttt{application.yml}, WS XML \textrightarrow{} API REST + OpenAPI,
        X.509 / SSL \textrightarrow{} configuration SSL declarative). Cette
        correspondance prouve que \textbf{l'architecture cible reste identique} :
        seule la mise en oeuvre gagne en concision, en securite et en
        maintenabilite.
\end{enumerate}

Ce choix presente un troisieme avantage, en phase directe avec les debouches du
parcours DSI. Le cahier des charges prevoit un \textbf{volet capteurs et analyse
predictive} (detection d'essaimage, traitement acoustique). Ce module se prete
naturellement a un \textbf{microservice Python} (pandas, scikit-learn) expose au
back-end Java via REST, illustrant concretement le metier d'\textbf{analyste de
donnees} cite au plan d'etudes et l'ouverture d'une \textbf{architecture
faiblement couplee} a l'heterogeneite technologique.

\begin{encadre}
\textbf{Conclusion :} le projet satisfait pleinement les exigences academiques
(J2EE) tout en documentant sa \textbf{trajectoire de modernisation} (Spring Boot
et microservice data). Cette double lecture transforme une contrainte
pedagogique en atout professionnel : elle demontre non seulement la maitrise des
fondamentaux du systeme d'information, mais aussi la capacite a faire evoluer une
architecture vers les standards recherches sur le marche de l'emploi en 2026.
\end{encadre}